\documentclass[11pt]{article}

\usepackage[utf8]{inputenc}
\usepackage[T1]{fontenc}
\usepackage[spanish]{babel}
\usepackage[margin=2.3cm]{geometry}
\usepackage{xcolor}
\usepackage{enumitem}
\usepackage{titlesec}
\usepackage{booktabs}

% --- Estilo ---
\definecolor{azul}{rgb}{0.10,0.35,0.55}
\definecolor{verde}{rgb}{0.10,0.45,0.40}
\definecolor{gris}{gray}{0.40}

\titleformat{\section}{\large\bfseries\color{azul}}{}{0pt}{}
\titleformat{\subsection}{\normalsize\bfseries\color{verde}}{}{0pt}{}

\newcommand{\slide}[1]{\textbf{\color{azul}[Diapositiva #1]}}
\newcommand{\transicion}[1]{\par\smallskip\noindent\textit{\color{gris}Transición:} #1}

\setlist[itemize]{itemsep=2pt, topsep=3pt, leftmargin=1.4em}

\title{\textbf{Robochess} --- Guión de la presentación}
\author{Juan Pedro García Sanz \and Pablo Revuelto de Miguel \and Adolfo Peña Marín}
\date{Robótica Inteligente --- Universidad Loyola Andalucía}

\begin{document}
\maketitle

% ============================================================
\section*{Cómo nos organizamos}
% ============================================================

\noindent
Tres bloques, uno por persona, siguiendo el hilo natural de la presentación:
un sistema que \textbf{ve, oye, piensa y mueve}.

\begin{center}
\begin{tabular}{@{}llc@{}}
\toprule
\textbf{Persona} & \textbf{Bloque} & \textbf{Diapositivas} \\
\midrule
Juan Pedro & Introducción y arquitectura        & 1--4  \\
Adolfo     & Percepción: visión, voz y cerebro  & 5--7  \\
Pablo      & Movimiento, integración y cierre   & 8--12 \\
\bottomrule
\end{tabular}
\end{center}

\noindent
\textbf{Duración objetivo:} 10--12 min en total ($\sim$3--4 min por persona).
Hablad despacio, una idea por diapositiva. No hay que leer la diapositiva:
ella tiene el titular, vosotros contáis el porqué.

\bigskip
% ============================================================
\section{Bloque 1 --- Juan Pedro (diapositivas 1--4)}
% ============================================================

\subsection{\slide{1} Portada}
\begin{itemize}
    \item Saludar y presentar al equipo brevemente.
    \item Frase gancho: \emph{``Os vamos a enseñar a dos brazos robóticos
          jugando al ajedrez entre ellos, y lo único que tenéis que hacer es
          decirles la jugada en voz alta''.}
\end{itemize}

\subsection{\slide{2} Índice}
\begin{itemize}
    \item En 10 segundos: las cuatro partes que vais a recorrer
          (arquitectura, percepción, decisión y movimiento, resultados).
\end{itemize}

\subsection{\slide{3} Qué es Robochess}
\begin{itemize}
    \item Explicar la idea: dos ABB IRB120 enfrentados, partida autónoma en
          simulación, jugada dictada por voz.
    \item Insistir en el reto real: lo difícil \textbf{no} es mover un brazo,
          es \textbf{integrar cuatro mundos} que normalmente van por separado
          ---visión, voz, razonamiento y planificación--- en un mismo flujo.
\end{itemize}

\subsection{\slide{4} Arquitectura}
\begin{itemize}
    \item Apoyarse en el diagrama: voz y cámara entran al cerebro, el cerebro
          ordena el movimiento, y la cámara verifica el resultado.
    \item Mencionar que todo corre sobre \textbf{ROS\,2 Jazzy} y que las partes
          se comunican por \textbf{topics} (flujo continuo), \textbf{servicios}
          (petición--respuesta) y \textbf{acciones} (tareas largas).
    \item Idea clave: arquitectura \textbf{modular}; cada bloque se desarrolla y
          se prueba por separado.
\end{itemize}

\transicion{``Empecemos por cómo el sistema percibe el mundo: la vista y el
oído. Adolfo os lo cuenta.''}

\bigskip
% ============================================================
\section{Bloque 2 --- Adolfo (diapositivas 5--7)}
% ============================================================

\subsection{\slide{5} Visión: del píxel a la casilla}
\begin{itemize}
    \item Cámara cenital + detector \textbf{YOLOv8}: localiza cada pieza y la
          clasifica en \textbf{12 clases} (6 tipos $\times$ 2 colores).
    \item Explicar la \textbf{homografía}: cómo pasamos del centro de la caja
          en píxeles a una casilla del tablero (a1--h8).
    \item Matiz importante: la visión \textbf{verifica, no manda}. Contrasta lo
          que ve con el estado interno; la autoridad es el motor de ajedrez.
\end{itemize}

\subsection{\slide{6} Voz: del audio a una jugada legal}
\begin{itemize}
    \item Recorrido: \textbf{VAD} (corta cuando dejas de hablar) $\rightarrow$
          \textbf{Whisper} (audio a texto) $\rightarrow$ \textbf{LLM} (texto a
          notación de ajedrez, p.\,ej. \texttt{e2e4}).
    \item La decisión que nos salvó: el parser \textbf{solo} puede devolver una
          jugada que esté en la lista de \textbf{movimientos legales}; nunca se
          inventa nada.
    \item Red de seguridad: si el LLM falla, un parser por \textbf{expresiones
          regulares} toma el relevo. Robusto ante transcripciones ruidosas.
\end{itemize}

\subsection{\slide{7} El cerebro: las reglas mandan}
\begin{itemize}
    \item El motor (\textbf{python-chess}) guarda el estado autoritativo en
          formato \textbf{FEN} y valida cada jugada: legalidad, turno, y detecta
          \textbf{jaque, mate y tablas}.
    \item Es el \textbf{director de orquesta}: si la jugada vale, la empaqueta y
          se la pasa al movimiento mediante una \textbf{acción} ROS\,2; al
          terminar, actualiza estado y cede el turno.
\end{itemize}

\transicion{``Ya tenemos una jugada válida. Ahora hay que ejecutarla con el
brazo sin tirar medio tablero. Te lo cuenta Pablo.''}

\bigskip
% ============================================================
\section{Bloque 3 --- Pablo (diapositivas 8--12)}
% ============================================================

\subsection{\slide{8} Movimiento: planificar sin chocar}
\begin{itemize}
    \item Cada brazo tiene \textbf{6 grados de libertad}. \textbf{MoveIt\,2}
          resuelve la \textbf{cinemática inversa} y busca con \textbf{OMPL} una
          trayectoria \textbf{libre de colisiones} (tablero, piezas, otro brazo).
    \item \textit{Pick \& place} por fases: aproximar, agarrar, elevar,
          transportar, soltar, retirarse. Agarre \textbf{magnético} para piezas
          pequeñas.
    \item Detalle de diseño que merece mención: reutilizamos un único
          planificador en marcha en vez de levantar otro, para no duplicar trabajo.
\end{itemize}

\subsection{\slide{9} El escenario simulado}
\begin{itemize}
    \item \textbf{Gazebo Harmonic}: mesa, tablero, los dos IRB120 y la cámara,
          con física realista de peso y contactos.
    \item Las articulaciones se controlan con \textbf{ros2\_control} por el
          puente ROS--Gazebo. Subrayar: es la \textbf{misma} infraestructura que
          se usaría sobre un robot real.
\end{itemize}

\subsection{\slide{10} Qué funciona y qué nos costó}
\begin{itemize}
    \item Repasar rápido lo que funciona: partida completa por voz, validación
          en tiempo real, pick \& place con los dos brazos, visión como validador.
    \item Ser honestos con los retos (esto da valor y suena a trabajo real):
          \textbf{capturar el micrófono bajo WSL2}, lograr que la \textbf{pieza
          siguiera al gripper} durante el transporte, y \textbf{sobrevivir a
          transcripciones ruidosas}.
\end{itemize}

\subsection{\slide{11} Conclusiones}
\begin{itemize}
    \item La idea fuerza: integrar IA (visión, voz, lenguaje) con robótica
          (planificación, control) sobre ROS\,2, y que todo funcione junto.
    \item La modularidad fue clave: cada pieza es sustituible.
    \item Próximos pasos: motor que juegue solo (Stockfish), salto a hardware
          real, y afinar visión y tiempos de planificación.
\end{itemize}

\subsection{\slide{12} ¿Preguntas?}
\begin{itemize}
    \item Cerrar agradeciendo y abrir el turno de preguntas.
\end{itemize}

\bigskip
% ============================================================
\section*{Turno de preguntas: quién responde qué}
% ============================================================
\begin{itemize}
    \item \textbf{Juan Pedro}: arquitectura, ROS\,2, integración general.
    \item \textbf{Adolfo}: visión (YOLO, homografía), voz (Whisper, LLM), reglas.
    \item \textbf{Pablo}: MoveIt, planificación, simulación, control.
\end{itemize}

\noindent
\textit{Si una pregunta pilla a alguien, mirad al compañero del área; no pasa
nada por repartir.}

\newpage
% ============================================================
\section*{Banco de preguntas probables (con respuesta técnica)}
% ============================================================

\noindent
El profesor avisó de que preguntará bastante. Aquí van las preguntas más
probables con una respuesta razonada. La idea no es soltarlas de memoria, sino
tener el argumento claro. \textbf{[JP]}~=~Juan Pedro, \textbf{[A]}~=~Adolfo,
\textbf{[P]}~=~Pablo (responde preferentemente quien lleva ese área).

% ---------- ROS 2 ----------
\subsection*{Sobre ROS\,2 e integración}

\textbf{1. ¿Por qué ROS\,2 y no un único programa monolítico? [JP]}
\begin{itemize}
    \item ROS\,2 nos da \textbf{desacoplamiento}: cada capacidad (visión, voz,
          motor, movimiento) es un nodo independiente que se desarrolla,
          lanza y depura por separado.
    \item Nos da ``gratis'' el \textbf{transporte} (DDS), el descubrimiento de
          nodos, el registro y herramientas de introspección
          (\texttt{ros2 topic echo}, \texttt{rqt}).
    \item Y nos permite \textbf{sustituir piezas} sin tocar el resto: por
          ejemplo cambiar Whisper por otro ASR sin que el motor se entere.
\end{itemize}

\textbf{2. ¿Cuándo usáis topic, servicio o acción, y por qué? [JP]}
\begin{itemize}
    \item \textbf{Topic} para flujo continuo sin respuesta: la imagen de la
          cámara (\texttt{/overhead\_camera/image}) y el estado del tablero.
    \item \textbf{Servicio} para petición--respuesta corta y síncrona:
          interpretar una frase (\texttt{/chess/voice/parse}).
    \item \textbf{Acción} para tareas largas con \textit{feedback} y posibilidad
          de cancelación: ejecutar la jugada (\texttt{/chess/execute\_move}),
          que tarda segundos y publica en qué fase va (\textit{approach, grasp,
          lift\dots}).
\end{itemize}

\textbf{3. ¿Cómo evitáis que el nodo se bloquee esperando una respuesta? [JP/P]}
\begin{itemize}
    \item Usamos \textbf{ejecutores multihilo} y lanzamos el pipeline de cada
          jugada en un \textbf{hilo aparte}. Si esperásemos la respuesta del
          servicio/acción dentro del \textit{callback}, el ejecutor se
          bloquearía y nunca llegaría esa respuesta (interbloqueo).
    \item En lugar de re-``spinear'' el nodo, esperamos el \textit{future} con un
          \textit{done-callback}, dejando el ejecutor libre para entregar el
          resultado.
\end{itemize}

\textbf{4. ¿Cómo se mueven realmente las articulaciones? [P]}
\begin{itemize}
    \item Con \textbf{ros2\_control}: un \texttt{joint\_trajectory\_controller}
          por brazo recibe la trayectoria que planifica MoveIt y la ejecuta a
          250\,Hz.
    \item En simulación, el plugin \texttt{gz\_ros2\_control} hace de
          \textit{hardware interface}; sobre un robot real se cambiaría ese
          plugin por el driver del ABB \textbf{sin tocar la lógica de arriba}.
\end{itemize}

% ---------- IA / integración ----------
\subsection*{Sobre la integración de la IA}

\textbf{5. ¿Cómo integráis modelos de IA pesados (Whisper, LLM) con ROS\,2? [A]}
\begin{itemize}
    \item Cada modelo vive en su \textbf{propio nodo} y se \textbf{carga de forma
          perezosa} (al arrancar el nodo, no al importar), para no penalizar el
          arranque del resto del sistema.
    \item Las dependencias de IA (torch, transformers, ultralytics, langchain)
          viven en un \textbf{entorno virtual} dedicado que se inyecta a los
          nodos por \texttt{PYTHONPATH}, manteniéndolas separadas del Python del
          sistema que usa ROS.
    \item Restricción real que respetamos: \textbf{numpy\,<\,2}, porque las
          extensiones C de ROS\,2 (\texttt{cv\_bridge}) están compiladas contra
          numpy\,1.x.
\end{itemize}

\textbf{6. ¿Dónde corre cada modelo: local o en la nube? [A]}
\begin{itemize}
    \item \textbf{Whisper}: en \textbf{local} (CPU), vía \texttt{transformers}.
          Pesos descargados una vez; sin coste ni API.
    \item \textbf{LLM} (Llama\,3): a través de un \textbf{endpoint de inferencia
          de HuggingFace} (en la nube). Por eso hay un \textit{fallback} local:
          si no hay red o token, el sistema sigue funcionando.
\end{itemize}

\textbf{7. ¿Cómo garantizáis que el sistema no haga una jugada inventada o
ilegal? [A]} \emph{(Pregunta estrella: la robustez es nuestro punto fuerte.)}
\begin{itemize}
    \item \textbf{Doble red}. Primero, al LLM le pasamos en el \textit{prompt}
          la \textbf{lista de movimientos legales} y le obligamos a devolver uno
          de ella o \texttt{null}; validamos el JSON de salida.
    \item Segundo, aunque el LLM acierte, el movimiento se \textbf{vuelve a
          validar} contra python-chess antes de mover el brazo. Si no es legal,
          no se ejecuta. El motor es la \textbf{autoridad}, no el LLM.
\end{itemize}

% ---------- LangChain / LangGraph ----------
\subsection*{Sobre LangChain / LangGraph}

\textbf{8. ¿Por qué LangChain y no LangGraph? [A]} \emph{(El profe lo va a
preguntar seguro.)}
\begin{itemize}
    \item LangGraph brilla en flujos \textbf{con estado, ramas, ciclos y
          múltiples agentes}. Nuestra tarea de parseo es de \textbf{un solo
          paso}: \textit{prompt} $\rightarrow$ LLM $\rightarrow$ JSON validado.
    \item Para eso, una \textbf{cadena} de LangChain (\texttt{prompt | llm |
          parser}) es la herramienta adecuada; LangGraph sería sobreingeniería.
    \item Dónde \textbf{sí} encajaría LangGraph: convertirlo en un agente con
          estado de partida, memoria y herramientas (consultar al motor,
          re-preguntar al usuario si la frase es ambigua). Es una vía de mejora
          honesta que podemos mencionar.
\end{itemize}

\textbf{9. ¿Qué es exactamente vuestra ``cadena'' de LangChain? [A]}
\begin{itemize}
    \item Tres eslabones con el operador \textit{pipe}: una
          \texttt{ChatPromptTemplate} (sistema + frase del usuario, con FEN y
          jugadas legales), el modelo de chat (\texttt{ChatHuggingFace} sobre un
          \texttt{HuggingFaceEndpoint}) y un \textit{parser} de salida.
    \item Usamos \textbf{temperatura 0} para que sea determinista: en parseo no
          queremos creatividad, queremos exactitud.
\end{itemize}

\textbf{10. ¿Por qué un LLM y no solo expresiones regulares? [A]}
\begin{itemize}
    \item El LLM entiende lenguaje \textbf{natural y variado} (``mueve el caballo
          a f3'', ``captura en d5 con la torre'', mezcla español/inglés) que un
          regex rígido no cubre.
    \item Pero mantenemos el \textbf{regex como fallback}: cubre los casos
          típicos sin red ni coste, y nos dio robustez cuando la transcripción
          venía con ruido. Es una decisión de \textbf{ingeniería}, no de moda.
\end{itemize}

% ---------- Visión ----------
\subsection*{Sobre el modelo de visión}

\textbf{11. ¿Qué modelo de visión usáis y por qué? [A]}
\begin{itemize}
    \item \textbf{YOLOv8} (Ultralytics): detector de objetos en \textbf{una
          pasada}, rápido y suficientemente preciso para 12 clases (6 piezas
          $\times$ 2 colores) desde una cámara cenital.
    \item Lo entrenamos para nuestras piezas; cargamos los pesos
          (\texttt{.pt}) por parámetro, así que el modelo es intercambiable.
\end{itemize}

\textbf{12. ¿Cómo pasáis de una detección en la imagen a una casilla? [A]}
\begin{itemize}
    \item Tomamos el \textbf{centro} de la caja detectada y aplicamos una
          \textbf{homografía} cámara cenital $\rightarrow$ plano del tablero,
          calculada a partir de la altura y el campo de visión de la cámara.
    \item Eso da coordenadas métricas en el tablero, que discretizamos a casilla
          (a1--h8) con la geometría del tablero (tamaño de casilla y centro).
    \item Para hardware real está previsto pasar a homografía con \textbf{4
          marcadores ArUco} en las esquinas.
\end{itemize}

\textbf{13. Si la visión se equivoca, ¿se rompe la partida? [A]}
\begin{itemize}
    \item No. La visión es \textbf{advisory}: solo \textbf{avisa} si lo que ve
          contradice el estado interno, y con \textbf{histéresis} (umbral de
          confianza, mínimo de detecciones y \textit{cooldown} para no spamear).
    \item Nunca modifica el estado ni bloquea la ejecución. La fuente de verdad
          es el motor de ajedrez.
\end{itemize}

% ---------- Voz ----------
\subsection*{Sobre el modelo de voz}

\textbf{14. ¿Por qué Whisper-small y no otro tamaño? [A]}
\begin{itemize}
    \item \textbf{Equilibrio} precisión/latencia en español sobre CPU.
          \textit{tiny}/\textit{base} son más rápidos pero fallan más;
          \textit{medium}/\textit{large} son más precisos pero lentos sin GPU.
    \item Es un \textbf{parámetro}: se puede subir de tamaño o mover a GPU
          (\texttt{device:=cuda}) sin tocar código.
\end{itemize}

\textbf{15. ¿Cómo sabéis cuándo empieza y acaba una orden de voz? [A]}
\begin{itemize}
    \item Con \textbf{detección de actividad de voz} por energía (RMS): un nodo
          de captura abre la grabación cuando la energía supera un umbral y la
          cierra tras $\sim$700\,ms de \textbf{silencio}.
    \item Así Whisper recibe \textbf{una intervención completa} y no audio
          cortado. (Detalle real: bajo WSL2 tuvimos que usar el backend
          PulseAudio porque el resto capturaba silencio.)
\end{itemize}

% ---------- Integración global ----------
\subsection*{Sobre la integración de todo}

\textbf{16. Contadme el viaje completo de ``mueve el peón de e2 a e4''. [todos]}
\begin{itemize}
    \item \textbf{[A]} El micrófono capta, la VAD segmenta, \textbf{Whisper}
          transcribe y el \textbf{LLM} traduce a \texttt{e2e4}, validado contra
          las jugadas legales.
    \item \textbf{[JP]} El \textbf{game manager} recibe la orden, la valida con
          python-chess y la despacha como \textbf{acción} al movimiento.
    \item \textbf{[P]} \textbf{MoveIt} planifica la trayectoria sin colisiones y
          \textbf{ros2\_control} la ejecuta; se hace el \textit{pick \& place} y,
          al terminar, el motor actualiza estado y cambia de turno.
\end{itemize}

\textbf{17. ¿Qué pasa con los dos brazos? ¿Se coordinan o chocan? [P]}
\begin{itemize}
    \item Hay \textbf{dos grupos de planificación} (\texttt{white\_arm} y
          \texttt{black\_arm}); cada jugada la ejecuta el brazo de su color.
    \item Como el juego es \textbf{por turnos}, no se mueven a la vez. Aun así,
          el planificador conoce el otro brazo como \textbf{obstáculo} en la
          escena, de modo que evita colisiones al planificar.
\end{itemize}

\textbf{18. ¿Qué es lo que más os costó integrar? [todos]}
\begin{itemize}
    \item La \textbf{frontera IA--ROS}: convivencia de versiones (numpy<2),
          entornos de Python y cargar modelos pesados sin bloquear el sistema.
    \item En manipulación, que la \textbf{pieza siguiera al gripper} de forma
          visualmente correcta durante el transporte.
    \item En voz, la \textbf{captura de micrófono} fiable y la robustez ante
          transcripciones ruidosas.
\end{itemize}

\end{document}
